\section{Conclusion, Limitations, and Future Directions}

In this paper, we introduce the (Nested) WAIT algorithm, which offers theoretical guarantees for online scheduling in large language model (LLM) inference problems, even when the output length is unknown upon arrival. Despite its theoretical properties and performance under idealized conditions, determining the optimal parameters for the (Nested) WAIT algorithm remains an unresolved challenge, as does the associated theoretical analysis. Furthermore, extending the algorithm to multi-GPU scenarios---prevalent in real-world applications---introduces additional complexities. Additionally, since more and more system-level optimization methods for LLM inference have been proposed, how to build on analytical frameworks for these newly prevalent models becomes necessary. In the future, we plan to strategically invest in research across the following directions.

\subsection{More Volatile Arrivals}

In practical settings, the distribution of incoming prompts fluctuates over time, violating the stationarity assumption underlying our model. While our algorithm can be directly adapted to scenarios with moderate variations in arrival rates—provided they remain within system capacity (as we show in Appendix \ref{sec:extension})—by re-solving the fluid equilibrium, a thorough analysis of its performance in more volatile environments, such as during demand surges, remains an open question. Investigating how the (Nested) WAIT algorithm behaves under such conditions is an important avenue for future research.

\subsection{Extension to Multi-GPU Systems}
Our current research is limited to single-GPU systems, capable of supporting LLMs such as Llama-13B. Extending the algorithm to multi-GPU environments introduces additional considerations, including communication costs, parallelization strategies, and hardware constraints. For example, while single-GPU deployment is relatively straightforward, multi-GPU setups necessitate careful management of data parallelism, tensor parallelism, and pipeline parallelism. Moreover, communication overhead becomes a critical factor, particularly since many large-scale models require multiple GPUs for deployment. Developing a theoretical framework that incorporates these elements poses a significant challenge, which we plan to investigate further.

\subsection{Analytical Framework for Up-to-Date LLM Inference Models}
In recent years, system-level designs for LLM inference have significantly influenced the field. For instance, the multi-head latent attention (MLA) mechanism introduced by \cite{deepseekai2024deepseekv2strongeconomicalefficient} has achieved a 90\% reduction in memory cost through an innovative compression technique. Integrating these state-of-the-art system-level optimizations into a theoretical framework presents a promising direction for future research.
